\documentclass[11pt,a4paper]{article}

\usepackage{fontspec}
\usepackage{polyglossia}
\setmainlanguage{russian}
\setotherlanguage{english}
\setmainfont{PT Serif}[Ligatures=TeX]
\setsansfont{PT Sans}[Ligatures=TeX]
\setmonofont{PT Mono}[Scale=0.92]
\newfontfamily\cyrillicfont{PT Serif}[Ligatures=TeX]
\newfontfamily\cyrillicfontsf{PT Sans}[Ligatures=TeX]
\newfontfamily\cyrillicfonttt{PT Mono}[Scale=0.92]

\usepackage{amsmath,amssymb}
\usepackage{booktabs}
\usepackage[margin=27mm,top=25mm,bottom=27mm]{geometry}
\usepackage{caption}
\usepackage{enumitem}
\usepackage{microtype}
\usepackage{tikz}
\usetikzlibrary{positioning,arrows.meta,decorations.pathreplacing,calligraphy}
\newcommand{\clip}{\operatorname{clip}}
\newcommand{\sgn}{\operatorname{sign}}

\captionsetup{font=small,labelfont=bf,labelsep=period}
\setlist{itemsep=3pt,parsep=0pt,topsep=5pt}
\setlength{\parskip}{0pt}
\setlength{\parindent}{1.2em}
\linespread{1.06}
\emergencystretch=2.2em

\title{\vspace{-16mm}\textbf{\Large Как мы предсказываем ингибирование цитохромов P450}\\[3mm]
\normalsize Устройство решения и математика за ним, без предварительных знаний}
\author{}
\date{\normalsize 9 сентября 2026}

\begin{document}
\maketitle
\thispagestyle{empty}

\begin{abstract}
\noindent
Мы участвуем в слепом соревновании OpenADMET: по структуре молекулы предсказать, насколько
сильно она подавляет четыре изофермента цитохрома P450. Этот текст объясняет, что мы построили
и почему, для читателя, который знает, что такое \textrm{IC}$_{50}$ и зачем нужны цитохромы,
но не занимается машинным обучением. Чисел здесь почти нет: интересны идеи. Формулы есть, но
все они собраны во врезках, каждый символ назван словами, и текст читается связно, если врезки
пропустить. Отдельный раздел посвящён тому, чего наше решение \emph{не} умеет --- это не
оговорка для приличия, а самая полезная часть.
\end{abstract}

\vspace{2mm}
\noindent\fbox{\parbox{\dimexpr\linewidth-2\fboxsep-2\fboxrule}{\small
\textbf{Как читать формулы.} Их пять, и все про одно и то же: как устроен счёт в
соревновании и что из этого следует. Обозначения сквозные --- $p$ всегда предсказание модели,
$y$ измеренное значение, $[\ell,h]$ полоса погрешности вокруг него. Формулы стоят в рамках и
дублируются словами рядом; если они мешают, их можно пропустить без потери смысла.}}

\vspace{3mm}
{\small\tableofcontents}
\vspace{4mm}

\section{Задача}

Цитохромы P450 метаболизируют большую часть лекарств, попадающих в организм. Если новое
соединение подавляет один из них, любой другой препарат, который этим ферментом
обезвреживается, начинает накапливаться. Так возникает лекарственное взаимодействие --- одна
из самых частых причин, по которым молекула вылетает из разработки на поздней стадии или
получает ограничения в инструкции.

Проверять это в эксперименте дорого, а кандидатов много. Поэтому есть смысл научиться отвечать
на вопрос \emph{до} эксперимента, по одной лишь структуре.

Организаторы дали обучающую выборку с измеренными значениями и спрятали набор соединений, для
которых надо прислать предсказания. Правильные ответы не покажут до конца соревнования,
попытка одна. Предсказать нужно две вещи: \textbf{силу обратимого ингибирования} для четырёх
изоферментов и \textbf{наличие ингибирования, зависящего от времени} --- того случая, когда
соединение не просто садится в активный центр, а выводит фермент из строя необратимо.

\section{Что нам дали}

Данные дырявые. Подавляющее большинство молекул измерено только на одном ферменте, на всех
четырёх --- считанные десятки. Это не дефект выборки, а обычное следствие того, как
накапливаются экспериментальные данные: разные соединения приходили из разных кампаний с
разными целями.

Кроме самих значений дали три вещи, и каждая оказалась важной.

\textbf{Границы доверия.} У каждого измерения есть не только точка, но и полоса, внутри
которой лежит истинное значение с учётом погрешности анализа. Для экспериментатора это
обычная вещь; как выяснится, это самая полезная часть данных.

\textbf{Одноточечный скрининг.} Отдельная таблица: для многих молекул есть не кривая, а одно
измерение при одной концентрации. Строк там заметно больше, чем полных кривых, но это другая
величина --- прямо в \textrm{IC}$_{50}$ она не переводится.

\textbf{Плечо преинкубации.} Те же молекулы, померенные после предварительной инкубации с
ферментом. Разница между этим значением и обычным --- и есть след необратимого ингибирования.

\section{Одна вещь, без которой дальше ничего не понятно}

Здесь надо сразу оговорить, кто что делает, потому что из этого следует всё остальное.

\textbf{Соревнование не делает с нашими числами ничего.} Оно берёт присланный файл и считает
штраф по нему как есть. А вот мы сами, последним шагом перед отправкой, прогоняем предсказания
через \textbf{перемасштабирование}: все предсказания одного фермента сдвигаем и стягиваем к
общему центру на одну и ту же величину. Это наше решение, а не правило соревнования. Но раз
оно принято, дальше всё следует из него.

\begin{center}\framebox{\parbox{0.93\linewidth}{\small
\textbf{Формула 1: перемасштабирование.} Каждое предсказание $p$ заменяется на
\[ q \;=\; c + \lambda\,(p - c), \qquad \lambda > 0, \]
где $c$ --- общий центр, а $\lambda$ --- общий коэффициент сжатия; оба подбираются один раз на
фермент. Словами: сдвинуть все предсказания и стянуть их к центру.

\smallskip
Ключевое в том, что $\lambda$ строго положительна. Если $p_1 > p_2$, то
$q_1 - q_2 = \lambda\,(p_1 - p_2) > 0$: \textbf{порядок сохраняется в точности}, какими бы ни
были $c$ и~$\lambda$. Это преобразование физически не способно переставить два соединения
местами.}}\end{center}

\subsection*{Что из этого следует}

\textbf{Любая систематическая ошибка у нас бесплатна.} Не потому, что метрика её прощает ---
она не прощает, и без этого шага модель была бы за неё наказана, --- а потому, что наш
собственный последний шаг её и так снимает. Считает модель все соединения активнее, чем они
есть, --- сдвиг уберёт разницу. Преувеличивает разброс --- сжатие уберёт и это.

Практический вывод неочевиден и стоил нам нескольких недель: \emph{улучшать калибровку модели
бессмысленно}. Всё, что можно поправить одним сдвигом и одним сжатием, поправит сдвиг и сжатие.

\textbf{А вот порядок этот шаг починить не может --- вообще никогда.} Если модель считает
соединение~А сильнее соединения~Б, а на самом деле наоборот, никакой сдвиг и никакое сжатие
этого не исправят: они монотонны.

\begin{quote}
Поэтому настоящая задача --- не «угадать \textrm{IC}$_{50}$», а \textbf{правильно расставить
соединения по силе}. Мы меряем себя ранговой корреляцией и относимся к ней как к главному
числу.
\end{quote}

Это звучит технической деталью, но это самый частый способ обмануться. Несколько раз мы
получали улучшение основной метрики, которое \emph{исчезало или меняло знак} после
перемасштабирования. Идея выглядела рабочей ровно до того момента, когда её мерили после того
шага, который всё равно будет применён.

\subsection*{Как и почему это работает у нас на самом деле}

Всё сказанное выше --- про то, что этот шаг \emph{способен} сделать. Полезно посмотреть, что
он делает \emph{фактически}, и ответ оказался не тем, которого мы ждали.

\textbf{Сжатия у нас практически нет.} Подобранный коэффициент $\lambda$ выходит равным единице
на трёх ферментах из четырёх и лишь немного меньше на четвёртом. Причина видна сразу, стоит
посмотреть на разброс: предсказания модели разбросаны в два-три раза \emph{у́же}, чем истинные
значения. Модель и так осторожна, сжимать нечего --- подгонка упирается в верхний край
допустимого диапазона и просит скорее растянуть, чем сжать.

\textbf{И на наших собственных данных этот шаг почти ничего не улучшает.} Если подобрать его
так, чтобы нашему же счёту было хорошо, выигрыш составляет тысячные доли --- внутри шума.
Историю он себе заработал раньше, на более грубых версиях модели, которым было что исправлять;
по мере того как модель становилась лучше, работы у этого шага оставалось всё меньше.

\textbf{Зачем он тогда стоит в подаче?} Потому что он там делает не то, что можно было
подумать. Мы настраиваем его \emph{не} под наши данные, а под предположение о спрятанных:
что они в среднем активнее нашей обучающей выборки. Поэтому предсказания сдвигаются вверх
заметно сильнее, чем выбрали бы наши собственные метки.

\begin{center}
\small
\begin{tabular}{lcc}
\toprule
 & подобрано под наши данные & \textbf{подаётся} \\
\midrule
сдвиг предсказаний вверх & едва заметный & заметно больший \\
счёт на \emph{нашей} выборке & лучший возможный & \textbf{хуже} \\
\bottomrule
\end{tabular}
\end{center}

То есть это \textbf{ставка, а не оптимизация}. Мы сознательно ухудшаем свой результат на
известных нам данных ради того, чтобы оказаться правыми на неизвестных. Ставка окупится, если
спрятанные соединения действительно активнее; если они устроены как наши --- мы просто отдали
часть счёта задаром. Насколько это рискованно, обсуждается в разделе~\ref{sec:limits}: это
самое крупное непроверяемое допущение во всём решении.

Что важно: \textbf{на порядок соединений эта ставка не влияет никак} --- она монотонна.
Проиграть на ней можно только в той части счёта, которая и так почти целиком определяется
не нами.

\section{Полоса вокруг измерения, и главная идея решения}

Теперь про то, что дало самый большой выигрыш.

Представьте, что вы просите оценить массу навески. Истинное значение известно, но весы дают
погрешность. Ответ, попавший внутрь этой погрешности, --- не ошибка: требовать большего
бессмысленно, прибор всё равно не различит.

Метрика соревнования устроена ровно так же. Если предсказание попадает \textbf{внутрь полосы
доверия} измерения, ошибка засчитывается нулевой. Штраф начинается только за выход за край,
и считается как расстояние до ближайшей границы, а не до точки.

\begin{center}\framebox{\parbox{0.93\linewidth}{\small
\textbf{Формула 2: во что нам обходится одно предсказание.} Пусть $p$ --- предсказание, а
$[\ell, h]$ --- нижняя и верхняя границы полосы измерения. Штраф за это соединение равен
\[ L(p) \;=\; \max\bigl(0,\; \ell - p,\; p - h\bigr). \]
Словами: если $p$ ниже нижней границы --- платим, насколько ниже; если выше верхней --- платим,
насколько выше; если внутри --- не платим ничего. Полный счёт по ферменту --- сумма таких
штрафов, делённая на такую же сумму для «глупого» предсказателя, который всем соединениям
выдаёт одно и то же среднее значение. Деление нужно, чтобы разные ферменты можно было
усреднить между собой; итог по четырём ферментам --- среднее четырёх таких дробей.}}\end{center}

Слева на рис.~\ref{fig:band} нарисована эта самая функция $L$. Обратите внимание на её форму:
плоское дно шириной в полосу погрешности и два прямых склона по краям.

\begin{figure}[t]
\centering
\begin{tikzpicture}[font=\small]
% ---------- левая панель: функция штрафа ----------
\begin{scope}
  \fill[black!8] (3,0) rectangle (5,2.3);
  \draw[->] (0.7,0) -- (7.7,0) node[right] {$p$};
  \draw[->] (0.7,0) -- (0.7,2.6) node[above] {$L(p)$};
  \draw[very thick] (1.0,2.0) -- node[sloped,above,inner sep=2pt] {$-1$} (3,0);
  \draw[very thick] (3,0)     -- node[above,inner sep=3pt] {$0$}          (5,0);
  \draw[very thick] (5,0)     -- node[sloped,above,inner sep=2pt] {$+1$}  (7,2);
  \draw[dashed,black!55] (3,0) -- (3,2.3);
  \draw[dashed,black!55] (5,0) -- (5,2.3);
  \draw[decorate,decoration={calligraphic brace,amplitude=4pt,mirror}]
       (3,-0.5) -- node[below,yshift=-3pt] {полоса погрешности} (5,-0.5);
  \node[below] at (3,-0.05) {$\ell$};
  \node[below] at (5,-0.05) {$h$};
  \node[align=center] at (4,3.05) {\textbf{штраф}};
\end{scope}
% ---------- правая панель: куда переносится цель ----------
\begin{scope}[xshift=9.6cm,yshift=0.55cm]
  \fill[black!8] (1.3,-0.3) rectangle (3.3,1.35);
  \draw[thick] (0.1,0) -- (4.9,0) node[right] {$p$};
  \draw[dashed,black!55] (1.3,-0.3) -- (1.3,1.35);
  \draw[dashed,black!55] (3.3,-0.3) -- (3.3,1.35);
  \node[below] at (1.3,-0.32) {$\ell$};
  \node[below] at (3.3,-0.32) {$h$};
  \fill (4.3,0.8) circle (2.3pt) node[above,inner sep=3pt] {$p_0$};
  \draw[->,thick] (4.22,0.8) -- (3.36,0.8);
  \node[below,inner sep=4pt] at (3.83,0.78) {цель};
  \fill (2.3,0.8) circle (2.3pt) node[above,inner sep=3pt] {$p_0'$};
  \node[align=center] at (2.5,2.3) {\textbf{мишень обучения}};
\end{scope}
\end{tikzpicture}
\vspace{7mm}
\caption{\textbf{Слева:} штраф за одно соединение как функция предсказания $p$; числа над
отрезками --- наклон. Дно плоское: внутри полосы погрешности метрика безразлична к тому, где
именно оказалось предсказание, и наклон там равен нулю. Снаружи наклон равен $\pm1$ независимо
от того, насколько далеко ушло предсказание, --- это и есть формула~4.
\textbf{Справа:} куда мы переносим цель обучения. Предсказание $p_0$, оказавшееся снаружи,
притягивается к ближайшему краю полосы, а не к точке измерения. Предсказание $p_0'$, уже
попавшее внутрь полосы, остаётся на месте: цель для него совпадает с ним самим, и переучивать
модель на нём нечему.}
\label{fig:band}
\end{figure}

Обычная модель этого не знает: её учат целиться в точку. Мы сделали иначе. Модель сначала
обучается как обычно, а потом \textbf{переобучается на исправленную цель}: если её предсказание
и так внутри полосы --- цель совпадает с предсказанием, менять нечего; если снаружи ---
целью становится ближайший край полосы, а не далёкая точка. Это правая панель рис.~\ref{fig:band}.

Смысл в том, что модель перестаёт тратить силы на то, за что её всё равно не штрафуют, и
занимается только теми соединениями, где действительно промахнулась.

\begin{center}\framebox{\parbox{0.93\linewidth}{\small
\textbf{Формула 3: почему это законно, а не просто правдоподобно.} Обозначим через
$\clip(p,\ell,h)$ операцию «загнать $p$ в полосу»: вернуть $\ell$, если $p$ ниже неё, вернуть
$h$, если выше, и вернуть само $p$, если внутри. Новая цель для соединения --- это
$t = \clip(p_0,\ell,h)$, где $p_0$ --- честное предсказание модели, полученное на данных,
которых она не видела при обучении.

\smallskip
Тогда для \emph{любого} нового предсказания $p$ верно
\[ \bigl|\,p - t\,\bigr| \;\ge\; L(p), \]
причём в точке $p = p_0$ здесь стоит равенство. Словами: расстояние до новой цели никогда не
меньше настоящего штрафа и совпадает с ним там, где модель сейчас находится. Значит, уменьшая
расстояние до цели, мы гарантированно уменьшаем и сам штраф --- мы давим на потолок, который
лежит поверх настоящей функции потерь и касается её в текущей точке. Приём известен как
мажорирование; он же стоит за многими алгоритмами вида «сделай шаг по удобной задаче, повтори».}}\end{center}

У той же картинки есть второе, более короткое прочтение --- через наклон.

\begin{center}\framebox{\parbox{0.93\linewidth}{\small
\textbf{Формула 4: наклон штрафа.} Производная функции штрафа по предсказанию равна
\[ \frac{dL}{dp} \;=\; \sgn\bigl(p - \clip(p,\ell,h)\bigr). \]
Проверяется перебором трёх случаев и видна прямо на левой панели рис.~\ref{fig:band}:
слева от полосы наклон равен $-1$, внутри полосы $0$, справа $+1$. Иначе говоря,
\textbf{направление, в котором надо двигать предсказание, полностью определяется тем, с какой
стороны от полосы оно оказалось} --- и ничем больше. Ни расстояние до точки измерения, ни
ширина полосы на это направление не влияют.

\smallskip
Отсюда и берётся вся конструкция: величина шага задаётся мажорированием из формулы~3,
а направление --- этим наклоном. Обе формулы проверены численно.}}\end{center}

\textbf{Это самый крупный вклад во всём решении} --- примерно треть накопленного улучшения
приносит одна эта идея. И она не про алгоритм, а про то, чтобы внимательно прочитать условие
задачи.

\section{Пять читателей одной молекулы}

Молекулу сначала переводят в числа: какие структурные фрагменты в ней встречаются, обычные
физико-химические дескрипторы (масса, липофильность, доноры водородной связи) и небольшой
набор признаков, написанных специально под цитохромы --- про основные атомы азота, способные
координировать гем, про размер, форму и заряд.

Затем пять разных моделей независимо читают эти числа, и каждая выдаёт своё предсказание.
Они устроены по-разному нарочно: одна учится на каждом ферменте отдельно; вторая --- на всех
сразу, с пометкой, о каком ферменте речь; третья опирается на сходство структур по принципу
«похожие молекулы ведут себя похоже»; четвёртая --- простая линейная, только на
физико-химических дескрипторах; пятая --- небольшая нейросеть с общим стволом и четырьмя
выходами.

Их предсказания усредняются. Идея та же, что за консилиумом: пять специалистов, ошибающихся
\emph{по-разному}, вместе надёжнее любого одного, даже если каждый по отдельности слабее
лучшего. Работает это только пока они ошибаются по-разному --- если все смотрят на одно и то
же, усреднение не даёт ничего. Ниже будет пример, где именно так и вышло.

\begin{center}
\begin{tikzpicture}[
  box/.style={draw,rounded corners=2pt,align=center,font=\small,inner sep=4pt,minimum height=8mm},
  wide/.style={box,minimum width=30mm},
  ar/.style={-{Latex[length=2mm]},thick}
]
\node[wide] (smiles) {структура};
\node[wide,right=7mm of smiles] (feat) {числовое\\описание};
\node[box,right=8mm of feat,minimum width=24mm] (m3) {пять моделей};
\node[wide,right=8mm of m3] (dz) {правка\\на «мёртвую зону»};
\node[wide,below=7mm of dz] (avg) {усреднение};
\node[wide,left=8mm of avg] (cal) {перемасштабирование};
\draw[ar] (smiles) -- (feat);
\draw[ar] (feat) -- (m3);
\draw[ar] (m3) -- (dz);
\draw[ar] (dz) -- (avg);
\draw[ar] (avg) -- (cal);
\end{tikzpicture}
\end{center}

Правка на «мёртвую зону» применяется к каждому участнику отдельно, до усреднения.

\section{Ферменты учат друг друга}

Одно наблюдение стоит отдельного раздела, потому что оно про биологию, а не про технику.

Модель, которая учится на всех четырёх ферментах сразу, предсказывает лучше, чем четыре
отдельные модели. Естественное объяснение --- «просто больше данных» --- оказалось неверным:
мы добавили заметно больше строк, максимально похожих на имеющиеся, и это не дало ничего.

Верное объяснение другое. Дело не в количестве, а в \textbf{контрасте}. Когда модель видит,
что одно и то же соединение сильно подавляет один фермент и не трогает другой, она узнаёт не
про соединение, а про то, \emph{чем эти ферменты отличаются}. Мы проверили это прямо: если
спрятать от модели, о каком ферменте идёт речь, она становится заметно хуже четырёх отдельных
моделей. А если сообщить ей средний уровень активности каждого фермента, но по-прежнему не
давать различать их в остальном, --- не помогает совсем.

Значит выигрыш даёт именно способность выучить \emph{разные} закономерности для разных
ферментов, а не знание про их общий уровень. Это, по сути, утверждение о том, что различия
между изоформами информативнее их сходства.

\begin{center}\framebox{\parbox{0.93\linewidth}{\small
\textbf{Формула 5: две гипотезы, которые легко перепутать.} Пусть $x$ --- описание молекулы,
$e$ --- номер фермента. Есть два способа устроить общую модель:
\[
\text{уровень:}\quad p = f(x) + \mu_e
\qquad\text{против}\qquad
\text{контраст:}\quad p = f_e(x).
\]
Слева одна и та же зависимость от структуры, сдвинутая на свою константу $\mu_e$ для каждого
фермента. Справа --- своя зависимость на каждый фермент.

\smallskip
Эксперимент различает их напрямую. Если убрать $e$ совсем, оставив $p=f(x)$, модель становится
заметно хуже четырёх отдельных. Если вернуть только $\mu_e$, то есть выдать уровень каждого
фермента даром, --- не помогает совсем. Работает лишь третий вариант. \textbf{Платит не сдвиг
$\mu_e$, а зависимость $f_e$}.}}\end{center}

\begin{figure}[t]
\centering
\begin{tikzpicture}[scale=1.0]
\begin{scope}
  \draw[->,thick] (0,0) -- (4.6,0) node[right,font=\small] {структура};
  \draw[->,thick] (0,0) -- (0,3.0) node[above,font=\small] {активность};
  \foreach \d/\c in {0.35/black!35, 0.95/black!55, 1.55/black!75, 2.15/black}
    \draw[very thick,\c] (0.3,\d) -- (4.2,\d+0.62);
  \node[font=\small,align=center] at (2.3,-0.75) {\textbf{уровень}\\одна зависимость,\\четыре сдвига};
\end{scope}
\begin{scope}[xshift=7.4cm]
  \draw[->,thick] (0,0) -- (4.6,0) node[right,font=\small] {структура};
  \draw[->,thick] (0,0) -- (0,3.0) node[above,font=\small] {активность};
  \draw[very thick,black!35] (0.3,0.35) -- (4.2,2.55);
  \draw[very thick,black!55] (0.3,1.15) -- (4.2,1.45);
  \draw[very thick,black!75] (0.3,2.45) -- (4.2,0.55);
  \draw[very thick,black]    (0.3,1.85) -- (4.2,2.75);
  \node[font=\small,align=center] at (2.3,-0.75) {\textbf{контраст}\\своя зависимость\\на каждый фермент};
\end{scope}
\end{tikzpicture}
\vspace{5mm}
\caption{Две гипотезы о том, что общая модель извлекает из четырёх ферментов сразу. Слева ---
ферменты отличаются только уровнем; справа --- они по-разному отзываются на одни и те же черты
структуры, вплоть до противоположного знака. Измерение выбрало правую картину: выдать модели
уровень каждого фермента даром не помогает вовсе, а вот запретить ей различать ферменты ---
хуже, чем вообще не объединять данные.}
\label{fig:contrast}
\end{figure}

\section{Вклад, измеренный однажды, --- не вклад сегодня}

У предыдущего раздела есть продолжение, и оно оказалось важнее самого раздела.

Все элементы решения добавлялись по одному, и вклад каждого мерился в момент добавления:
насколько лучше стало от того, что мы его внесли. Так накопился список: вот эта идея стоит
столько-то, вот эта --- столько-то. Список сходился с общим результатом, и это выглядело
подтверждением.

Недавно мы проверили тот же список \emph{наоборот}: не «сколько добавляет элемент к тому, что
было», а «сколько теряется, если убрать его из готовой модели». Разница оказалась большой.

Кое-что подтвердилось почти точь-в-точь --- правка на «мёртвую зону» и нейросеть стоят ровно
столько, сколько записано. Зато та самая модель, которая учится на всех ферментах сразу,
в готовом ансамбле \textbf{не добавляет уже ничего}: без неё получается даже чуть-чуть лучше.
Механизм с контрастом реален, но к моменту, когда усреднены четыре остальных, всё, что этот
участник приносил, там уже есть.

Это общее свойство, а не курьёз. Пока система растёт, элементы начинают дублировать друг
друга, и вклад, измеренный на пустом месте, перестаёт что-либо говорить о вкладе здесь и
сейчас. Единственный способ узнать правду --- выбивать элемент из готовой системы, а не
добавлять его к пустой.

\section{Как мы себя проверяем}

Эта часть ближе всего к тому, как устроена работа в лаборатории, и мы относимся к ней
серьёзнее, чем к самим моделям.

\textbf{Проверка на неродственниках.} Легко получить красивое число, если проверять модель на
молекулах, у которых в обучении лежат близкие аналоги: тогда достаточно узнать знакомое.
Мы разбиваем данные так, чтобы структурно похожие соединения попадали в одну группу целиком,
и группа целиком уходит либо в обучение, либо в проверку.

\textbf{Шумовой порог.} Прежде чем радоваться улучшению, надо знать, какого размера бывает
случайная разница. Мы это померили: одна только перетасовка разбиения двигает результат
сильнее, чем типичный эффект, который хочется объявить улучшением. Всё, что меньше порога,
у нас не считается результатом --- даже если очень хочется.

\textbf{Условия --- до прогона.} Прежде чем запускать проверку идеи, мы записываем, какой
результат будем считать успехом. Это ровно то же, что регистрация протокола до эксперимента,
и по той же причине: иначе критерий незаметно подстраивается под то, что получилось.

\textbf{Журнал.} Всё измеренное лежит в одном файле, устроенном как журнал \emph{неудач}:
верная идея занимает один пункт, неверная --- три (выдвижение, опровержение, поправка к
опровержению). Соотношение при чтении выходит примерно три к одному в пользу провалов, и это
честная картина того, как идёт работа.

\section{Что мы про себя знаем плохого}
\label{sec:limits}

\textbf{На той метрике, которую показывает турнирная таблица, весь наш выигрыш меньше шума
одного замера.} Разброс результата от одного прогона к другому больше, чем всё, что мы набрали
за месяцы. По рангу мы уверенно вне шума, но ранг в таблице не показывают. Это надо знать
заранее, чтобы не делать выводов из одной цифры на промежуточном этапе.

\textbf{Наша проверка строже настоящего экзамена, и не в ту сторону.} Мы намеренно проверяем
себя на молекулах без близких родственников. А спрятанные соединения устроены иначе: они
ближе к обучающей выборке, чем обучающая выборка сама к себе. Грубо говоря, у каждого пятого
тестового соединения есть близкий аналог среди наших --- а в нашей проверке такое встречается
примерно в одном случае из ста. Мы измеряем себя в режиме, в котором экзамена почти не будет.

\textbf{Один из четырёх ферментов ведёт себя иначе, и мы не до конца понимаем почему.}
На CYP3A4 почти всё, что помогает остальным, помогает слабее или не помогает вовсе. В итоге
для него подаётся не коллегия, а один-единственный её участник. Отдельно выяснилось, что
обучающая выборка по этому ферменту --- это на самом деле два разных набора, склеенных вместе:
широкий скрининг и аналоговая серия.

\textbf{Мы не знаем, насколько тестовый набор смещён по активности} относительно обучающего.
Финальное перемасштабирование настраивается под предположение об этом смещении, и это
предположение --- самое крупное допущение во всём решении.

\section{Второй трек: обратимое против необратимого}

Здесь получился результат в двух частях. Первая --- про то, как правильно поставить задачу;
вторая --- про то, что в самих данных нашлась поломка.

\subsection*{Задачу поставили неверно, и это видно из условия}

Надо предсказывать метку «зависит от времени» --- да или нет. Мы долго строили под неё
классификаторы и упирались в потолок. Оказалось, зря: \textbf{организаторы задали эту метку
формулой, и формула напечатана в условии}. Соединение помечается положительным, если после
преинкубации оно подавляет фермент как минимум вдвое сильнее, чем без неё, --- плюс оговорка про
слабые соединения, у которых сдвиг нельзя измерить напрямую. Никакой скрытой закономерности тут
нет; это определение, а не наша находка.

Находка в другом --- в том, что из этого следует. Раз метка \emph{вычисляется} из двух измеренных
величин (потентности в обычном плече и в плече с преинкубацией), учить «что такое TDI» не нужно
вовсе. Достаточно предсказать эти две физические величины --- ровно те же, что и на первом
треке, --- и подставить их в готовую формулу. Мы так и делаем. На соединениях, где оба плеча
измерены, формула воспроизводит опубликованную метку без единого исключения (2334 из 2334 для
CYP3A4, 1493 из 1493 для CYP2D6) --- это не заслуга модели, а подтверждение, что мы правильно
прочли условие.

Из той же формулы читается, почему два фермента даются по-разному. Потентность --- необходимое
условие: слабое соединение не может быть положительным, как бы ни сдвигалось. На CYP3A4 это
условие отсекает около 40\,\% соединений, и распознать их --- отдельная задача про силу, которую
модель решает хорошо. На CYP2D6 оно отсекает лишь около 6\,\%, поэтому там метка --- это почти
целиком вопрос сдвига. Вся разница в качестве между двумя ферментами объясняется этим одним
числом, и никакого другого объяснения не требует.

\subsection*{Тысяча меток оказалась заглушками}

Проверяя формулу, мы наткнулись на то, что стоит рассказать отдельно, потому что это про
целостность данных, а не про модель.

У части соединений измерено только плечо с преинкубацией, а обычного плеча нет. По CYP3A4 таких
\textbf{1249}, и \emph{все до одного} помечены отрицательными --- при том, что в среднем по этому
ферменту положителен каждый третий. Более того: у 1055 из этих 1249 (около 85\,\%) потентность
после преинкубации выше порога, за которым только и возможен положительный ответ, и по потентности
эти соединения \emph{сильнее} размеченной популяции, а не слабее (медиана 5.4 против 4.6).

Иными словами, это ровно те соединения, среди которых положительных должно быть особенно много,
а стоит ноль. Если бы метки были настоящими измерениями, хоть сколько-то из 1249 оказались бы
положительными; их ноль. Единственное разумное прочтение --- \textbf{это не измерения, а
заполнитель}: пустые клетки, залитые словом «нет». Практический вывод прямой: эти строки нельзя
давать классификатору, иначе он учится на тысяче перевёрнутых меток. Мы их исключаем.

\section{Итог}

Наше решение --- это не одна большая модель, а пять средних, усреднённых, с двумя правками
до и после: одна учит модель целиться в полосу измерения, а не в точку, вторая приводит
итоговую шкалу в порядок.

Больше всего выигрыша принесли не алгоритмы, а внимательное чтение условия: то, что метрика
прощает попадание в полосу погрешности. Второе по важности --- собственное наблюдение, а не
правило: общий сдвиг можно снять одним дешёвым шагом в самом конце, а значит бороться с ним
внутри модели незачем, и всё внимание надо отдать порядку.

И, пожалуй, главное, что мы вынесли. Из идей, которые кажутся разумными, работает примерно
одна из четырёх. А из тех, что уже сработали когда-то, часть перестаёт работать, когда вокруг
них вырастает всё остальное, --- и заметить это можно, только если время от времени разбирать
собственную конструкцию обратно.

\end{document}
